\documentclass[11pt]{article}
\usepackage[margin=1in]{geometry}
\usepackage{amsmath,amssymb,amsthm}
\usepackage{mathtools}
\usepackage{enumitem}
\usepackage{hyperref}
\usepackage{cleveref}
\usepackage{booktabs}
\usepackage{tabularx}
\usepackage{microtype}
\usepackage{cite}
\usepackage{xurl}
\usepackage{path}
\hypersetup{colorlinks=true,linkcolor=blue,citecolor=blue,urlcolor=blue}
\newtheorem{definition}{Definition}
\newtheorem{lemma}{Lemma}
\newtheorem{remark}{Remark}
\title{Public Request Notices for Source-Only Releases\\\large Outward-Facing Sentences for Within-Release Status and Cross-Release Carryforward}
\author{Anonymous}
\date{Unified archive note v0.13 (2026-03-28; archive v740)}
\begin{document}
\maketitle

\begin{abstract}
The archive can now say exactly what a source-only paper promises on request, which omitted-support families that promise covers, when those families are complete, what the within-release paper-level public token is, and whether that token survives or reopens under a successor release. One smaller publication-layer question still recurs in terse papers and release notes: what exact short sentence should readers actually see? This note gives that wrapper one home. It defines compact \emph{public request notices}: outward-facing sentence-and-pointer objects that compress either a within-release public request status envelope or a successor-release public request carryforward envelope without re-owning the narrower status, completion, or release semantics they summarize.
\end{abstract}

\paragraph{Non-redundancy note.}
Synthesis~56 owns exact paper-facing public request contracts, Synthesis~58 owns family-scoped completion criteria, Synthesis~59 owns the within-release paper-level public status token, Synthesis~60 owns the cross-release carryforward rule, Synthesis~62 owns the canonical sentence family and slot realization rule for source-only notices, Synthesis~63 owns the token-keyed choice of which such family applies, Synthesis~73 owns the request-side terminal stop later terse papers should usually cite, Synthesis~74--75 own the paired terminal / cutover discipline, and Synthesis~17 owns the concrete worked instantiation.\cite{series:publicrequestcontracts,series:requestfulfillmentcertificates,series:publicrequeststatusenvelopes,series:publicrequestcarryforward,series:publicrequestnoticenormalforms,series:publicrequestnoticeselections,series:publicrequestresponsepacketrefreshresponsepacketclosureverdicts,series:pairedterminalcitationdiscipline,series:pairedterminalcutoverrule,series:workedexample}
This note owns only the public wrapper sentence and reader-pointer order for those already-owned source-only status objects. Even when the same imported status or carryforward token names recur across many visible wrappers, those tokens remain basis vocabulary rather than notice ownership, and later terse papers should cite this note directly only when the exact outward-facing sentence or pointer order is itself live.

\paragraph{Scope.}
This is not a new public request contract, not a new completion rule, not a new paper-level status token, and not a new carryforward verdict.
It is a publication-interface note for the short sentence a reader actually sees once those narrower owners already exist.

\paragraph{Exports.}
This note exports (i) public request notices, (ii) a basis-first rule, (iii) a no-new-semantics compression rule, (iv) a wrapper-agreement rule, (v) a public request-notice sentence normal form saying that later terse papers may render one already-owned visible wrapper as one exact clause --- imported basis key, selected normal-form key, emitted notice key, and first reader pointer --- instead of restitching that clause from the forced matrix or widening immediately to the family / selector owners, and (vi) a minimal citation boundary separating the outward-facing source-only sentence from the status or carryforward owner that makes it true. If later notes instead need the canonical sentence family and slot realization rule behind that emitted sentence, they should import Synthesis~62 rather than widen this note. If they need the token-keyed choice of which such family applies, they should import Synthesis~63 rather than widen this note. If they need only the stable request-side stop for the adjacent visible source-only branch, they should cite Synthesis~73. If they need one concrete checked-answer/request-tail seam, they should drop to Synthesis~17. And if they need only the abstract paired-tail discipline or abstract stop-versus-widen judgment once that request-side stop is already fixed, they should cite Synthesis~74--75 rather than widen this wrapper note.

\paragraph{Later-terse-paper citation posture.}
Later terse papers should cite Synthesis~61 only when the live issue is the exact outward-facing source-only sentence or reader-pointer order that readers were shown. Stop at Synthesis~59 when the live issue is which within-release paper-level public token is true. Stop at Synthesis~60 when the live issue is whether that token survived, advanced, or reopened across one successor release. Stop at Synthesis~62 when the live issue is the reusable canonical notice family or slot map behind that visible sentence, and stop at Synthesis~63 when the live issue is the token-keyed choice of which such family applies. Widen through Synthesis~64--72 only when the visible wrapper is already fixed and the live issue has moved into the exact follow-up handoff, predecessor packet cut, successor reuse, refresh, wrapper, or successor-facing packet chain. Widen to Synthesis~73 when the live issue is whether that visible request-side branch already closes. Drop to Synthesis~17 when one concrete checked-answer/request-tail seam is the live issue instead. Widen to Synthesis~74--75 only when the remaining issue is the abstract paired-terminal discipline or abstract fail-closed cutover.\cite{series:publicrequeststatusenvelopes,series:publicrequestcarryforward,series:publicrequestnoticenormalforms,series:publicrequestnoticeselections,series:publicrequestresponsemenus,series:publicrequestresponsepackets,series:publicrequestresponsepacketcarryforward,series:publicrequestresponsepacketdeltaledgers,series:publicrequestresponsepacketrefreshnotices,series:publicrequestresponsepacketrefreshnoticenormalforms,series:publicrequestresponsepacketrefreshnoticeselections,series:publicrequestresponsepacketrefreshresponsemenus,series:publicrequestresponsepacketrefreshresponsepackets,series:publicrequestresponsepacketrefreshresponsepacketclosureverdicts,series:pairedterminalcitationdiscipline,series:pairedterminalcutoverrule,series:workedexample}


\paragraph{One five-stop visible-request route.}
Once a source-only branch has already fixed the imported within-release or successor-facing basis object, later terse papers should usually traverse one recurring five-stop midbody route: Synthesis~61 for the exact emitted notice and first reader pointer, Synthesis~62 for the reusable canonical sentence family behind that notice, Synthesis~63 for the token-keyed choice of which such family applies, Synthesis~64 for the smallest-sufficient follow-up handoff once that choice is fixed, and Synthesis~65 for the exact predecessor-side bounded packet cut that should actually travel. This note is only the first stop in that route. Stop here when the live issue is the exact visible wrapper clause readers saw. Hand off to Synthesis~62--65 as soon as the live issue becomes the reusable family, the selector, the handoff owner, or the exact carried slice rather than the emitted notice itself.

\paragraph{A tiny five-stop drill.}
In the worked successor request-scope branch, the honest split is: Synthesis~61 says which visible carryforward notice and first pointer readers saw; Synthesis~62 says which canonical family that sentence instantiates; Synthesis~63 says why that family and no other is licensed by the imported carryforward bundle; Synthesis~64 says that a \nolinkurl{request_scope_check} follow-up selects the contract owner; and Synthesis~65 says which exact bounded contract packet actually travels. A later terse paper that compresses all five into the emitted wrapper sentence has already overstated what the wrapper owner knows.

\section{Why notices deserve their own note}
A very short source-only paper often wants one sentence that is suitable for the visible paper body, for a review response, or for a successor release note.
By the time that sentence is written, the archive may already know the exact packet/menu cut, the family-scoped completion state, the within-release paper-level status token, and the cross-release carryforward verdict.
Without a separate wrapper owner, later notes tend to cite a sentence as if it owned the token itself.
That is the same compression mistake the archive already avoided for release-status lineage notices.
A public request notice keeps the prose short while leaving the public request status envelope and carryforward envelope as the semantic owners.

\begin{definition}[Public request notice]
Fix either a within-release public request status envelope $E$ or a successor-release public request carryforward envelope $C$.
A \emph{public request notice} is a tuple
\[
N=(\beta,\tau,\sigma,\Pi)
\]
consisting of one imported basis object $\beta$ (either a within-release public request status envelope or a successor-release public request carryforward envelope), one notice kind $\tau$ (either \texttt{within\_release} or \texttt{successor\_carryforward}), one outward-facing sentence $\sigma$, and one ordered reader-pointer list $\Pi$ naming the first maintained artifacts a reader should inspect if that sentence is challenged.
The notice owns the outward-facing sentence and pointer order.
It does not replace the imported basis object or the narrower contract/completion owners that basis object already summarizes.
\end{definition}

\begin{definition}[Basis-first rule]
A public request notice satisfies the \emph{basis-first rule} if its first reader pointer is the imported basis object itself, and if every later pointer moves leftward only to already-owned source-only request artifacts that the basis object depends on, such as the public request contract, request-fulfillment certificates, or concrete successor comparison context.
\end{definition}

\begin{definition}[No-new-semantics compression rule]
A public request notice satisfies the \emph{no-new-semantics compression rule} if every contested semantic piece of its sentence $\sigma$ is implied by a field already present in the imported basis object or by an explicit pointer in $\Pi$ to the narrower owner artifact that makes that phrase true.
In particular, the notice may not introduce new promised families, new completion claims, or new release-transition claims that are absent from the imported basis object and absent from the pointed owner artifacts.
\end{definition}

\begin{definition}[Wrapper-agreement rule]
A public request notice satisfies the \emph{wrapper-agreement rule} if it names one already-selected public-request-notice normal-form key, copies that normal form's rendered sentence exactly, and orders its reader pointers by the minimal basis-first audit chain compatible with its notice kind.
For a within-release notice, that chain begins at the imported public request status envelope and then widens only to the family-scoped request-fulfillment certificates and the public request contract.
For a successor-carryforward notice, that chain begins at the imported public request carryforward envelope and then widens only to the inherited status envelope, the successor-facing public request contract, and the family-scoped request-fulfillment certificates.
So the emitted wrapper is fixed by one imported basis object together with one already-selected canonical family rather than by ad hoc prose.
\end{definition}

\begin{definition}[Public request-notice sentence normal form]
A public request notice with imported basis key \texttt{B} and selected normal-form key \texttt{F} exports a \emph{public request-notice sentence normal form} if it supplies exactly one clause:
\begin{quote}\small
Imported basis \texttt{B} with selected normal form \texttt{F} emits notice \texttt{N} and first points to \texttt{P}.
\end{quote}
Here \texttt{N} is the emitted public-request-notice key and \texttt{P} is the first reader-pointer artifact named by the basis-first audit chain. The clause is only a rendering of the already-owned basis-to-family-to-wrapper relation. It does not replace the imported status / carryforward owner, the canonical family or selector owner, or the later request-side terminal stop.
\end{definition}

\begin{lemma}[Notice-first audit rule]
If a public request notice satisfies the basis-first rule, the no-new-semantics compression rule, and the wrapper-agreement rule, then any challenge to the visible sentence resolves first to the notice and then immediately to the imported basis object and selected normal form that fixed it, while any challenge to the actual source-only status or carryforward token still resolves to the basis object and its narrower pointed owners rather than to the sentence wrapper.
\end{lemma}

\begin{proof}
By definition, the notice owns only the outward-facing sentence and reader-pointer order.
The basis-first rule ensures that the first audit stop is always the imported status or carryforward owner rather than a later packet or prose wrapper.
The no-new-semantics compression rule forbids the notice from silently creating new promised families, completion verdicts, or cross-release claims.
The wrapper-agreement rule then pins the emitted sentence to one already-selected normal form and pins the first audit stops to the minimal basis-first chain compatible with the notice kind.
So a wording challenge can be audited through the notice, but the actual status or carryforward verdict remains owned by the imported basis object and the narrower owner artifacts it already points to.
That keeps the public sentence short without letting it become the semantic owner of the source-only request promise.\qedhere
\end{proof}

\begin{remark}[Minimal citation rule]
Cite the public request notice when the claim is about \emph{what outward-facing source-only sentence readers saw}.
Cite the public request status envelope when the claim is about \emph{which within-release paper-level public token is true}.
Cite the public request carryforward envelope when the claim is about \emph{whether that token survived, advanced, or reopened across a successor release}.
Cite the public request contract or request-fulfillment certificates directly when the claim is about exact promised families, delivered support paths, or completion details.
\end{remark}

\begin{remark}[Pointer stop discipline]
When later terse papers need only the stable request-side stop for the visible source-only branch, they should usually cite Synthesis~73, drop to Synthesis~17 for one concrete checked-answer/request-tail seam, and widen to Synthesis~74--75 only for the abstract stop-versus-widen judgment rather than expand this wrapper layer locally.
Widen through this note only when the exact outward-facing sentence or pointer order is itself disputed. In that case, for a within-release notice the first sufficient audit stop is the imported public request status envelope, and widening should proceed only to the family-scoped request-fulfillment certificates and then to the public request contract that fixes the promised cut.
For a successor-carryforward notice, the first sufficient audit stop is the imported carryforward envelope, and widening should proceed only to the imported within-release status envelope, the successor-facing public request contract, and then the family-scoped request-fulfillment certificates.
In both cases, repeated status-token and carryforward-token labels remain imported basis vocabulary rather than notice-owned semantics.
\end{remark}

\begin{table}[t]
\centering
\small
\begin{tabularx}{\linewidth}{@{}l l X@{}}
\toprule
Layer & Canonical object & Question answered \\
\midrule
Within-release basis & public request status envelope & Which paper-level source-only status token is true for this release? \\
Cross-release basis & public request carryforward envelope & Did that token survive, advance, or reopen under the successor cut? \\
Reader-facing wrapper & public request notice & What exact short source-only sentence and first pointers were shown to readers? \\
\bottomrule
\end{tabularx}
\caption{The source-only thread stays auditable when public sentences do not collapse into their status or carryforward owners.}
\label{tab:publicrequestnotices}
\end{table}

\paragraph{A forced public request-notice matrix.}
\begin{table}[t]
\centering
\small
\begin{tabularx}{0.97\linewidth}{@{}p{0.31\linewidth}X@{}}
\toprule
\textbf{Public row} & \textbf{Pinned value}\\
\midrule
Within-release forced wrapper & Basis key \nolinkurl{worked_example_receipt_interlock_public_request_status_envelope} together with selected normal form \nolinkurl{worked_example_receipt_interlock_within_release_complete_public_request_notice_normal_form} forces emitted notice \nolinkurl{worked_example_receipt_interlock_public_request_status_notice} \\
Successor forced wrapper & Basis key \nolinkurl{worked_example_receipt_interlock_public_request_carryforward_envelope} together with selected normal form \nolinkurl{worked_example_receipt_interlock_successor_preserved_complete_public_request_notice_normal_form} forces emitted notice \nolinkurl{worked_example_receipt_interlock_public_request_carryforward_notice} \\
Imported token floor & Both wrappers inherit \texttt{public\_request\_status\_token = complete} from their basis objects; the successor wrapper additionally inherits \texttt{carryforward\_token = preserved\_complete} as imported basis vocabulary rather than notice-owned semantics \\
Forced sentence floor & The within-release wrapper must copy the selected family's rendered sentence stating that the worked note's paper-level source-only request promise is \texttt{complete}; the successor wrapper must copy the selected family's rendered sentence stating that same complete status carries forward unchanged across the canned successor release \\
Forced pointer-order floor & Within release: status envelope $\to$ request-fulfillment certificates $\to$ public request contract. Successor carryforward: carryforward envelope $\to$ status envelope $\to$ public request contract $\to$ request-fulfillment certificates \\
Staleness trigger floor & If the imported basis tokens change, if a different normal form is selected, or if the first pointer ceases to name the imported basis object, this matrix can no longer be quoted unchanged because the emitted wrapper is no longer the same notice \\
Escalation pointer & Widen to Synthesis~59--60 for the token owners, Synthesis~62--63 for canonical-family / selector ownership, widen to Synthesis~73 for the request-side terminal stop, drop to Synthesis~17 for one concrete checked-answer/request-tail seam, and widen to Synthesis~74--75 only for the abstract paired-tail discipline or abstract cutover judgment \\
\bottomrule
\end{tabularx}
\caption{A forced public request-notice matrix. This is the smallest public answer later terse papers can quote directly when the live issue is the exact outward-facing source-only wrapper forced by one imported basis object plus one already-selected canonical family rather than the narrower token owner, the family / selector owners, or the request-side terminal tail.}
\label{tab:forced-public-request-notice-matrix}
\end{table}

This matrix is intentionally non-decorative: it pins the actual within-release and successor-facing basis keys, the selected normal-form keys, the forced emitted notice keys, and the forced first reader-pointer order without silently re-owning the within-release status token, the successor carryforward verdict, the canonical sentence family, or the request-side terminal stop.\cite{series:workedexample}

\paragraph{A tiny public request-notice sentence card.}
\begin{table}[t]
\centering
\small
\begin{tabularx}{0.97\linewidth}{@{}p{0.26\linewidth}X@{}}
\toprule
\textbf{Clause row} & \textbf{Exact visible clause}\\
\midrule
Within-release clause & Imported basis \nolinkurl{worked_example_receipt_interlock_public_request_status_envelope} with selected normal form \nolinkurl{worked_example_receipt_interlock_within_release_complete_public_request_notice_normal_form} emits notice \nolinkurl{worked_example_receipt_interlock_public_request_status_notice} and first points to \nolinkurl{example_public_request_status_envelope.json}. \\
Successor clause & Imported basis \nolinkurl{worked_example_receipt_interlock_public_request_carryforward_envelope} with selected normal form \nolinkurl{worked_example_receipt_interlock_successor_preserved_complete_public_request_notice_normal_form} emits notice \nolinkurl{worked_example_receipt_interlock_public_request_carryforward_notice} and first points to \nolinkurl{example_public_request_carryforward_envelope.json}. \\
Escalation pointer & Widen to Synthesis~59--60 for the token owners, Synthesis~62--63 for canonical-family / selector ownership, widen to Synthesis~73 for the request-side terminal stop, drop to Synthesis~17 for one concrete checked-answer/request-tail seam, and widen to Synthesis~74--75 only for the abstract paired-tail discipline or abstract cutover judgment. \\
\bottomrule
\end{tabularx}
\caption{A tiny public request-notice sentence card. This is the smallest public answer later terse papers can quote directly when the live issue is one exact outward-facing source-only wrapper clause rather than the narrower token owner, the family / selector owners, or the request-side terminal tail.}
\label{tab:public-request-notice-sentence-card}
\end{table}

This card is intentionally non-decorative: it renders the already-owned wrapper as one exact clause --- basis key, selected family key, emitted notice key, and first reader pointer --- without silently re-owning the within-release status token, the successor carryforward verdict, the canonical family, or the request-side terminal stop.\cite{series:workedexample}

\paragraph{One tiny wrapper-clause fill.}
For the worked within-release branch, fill the four clause slots with basis key \nolinkurl{worked_example_receipt_interlock_public_request_status_envelope}, normal-form key \nolinkurl{worked_example_receipt_interlock_within_release_complete_public_request_notice_normal_form}, notice key \nolinkurl{worked_example_receipt_interlock_public_request_status_notice}, and first pointer \nolinkurl{example_public_request_status_envelope.json}. For the worked successor branch, fill them with basis key \nolinkurl{worked_example_receipt_interlock_public_request_carryforward_envelope}, normal-form key \nolinkurl{worked_example_receipt_interlock_successor_preserved_complete_public_request_notice_normal_form}, notice key \nolinkurl{worked_example_receipt_interlock_public_request_carryforward_notice}, and first pointer \nolinkurl{example_public_request_carryforward_envelope.json}. Later terse papers may quote either one-clause wrapper only while those four anchors stay fixed; if any anchor drifts, reopen the narrower owner instead of repairing the visible clause ad hoc.\cite{series:workedexample}

\paragraph{One tiny wrapper-forcing drill.}
In the worked within-release wrapper, the imported basis token \texttt{public\_request\_status\_token = complete} together with selected family \nolinkurl{worked_example_receipt_interlock_within_release_complete_public_request_notice_normal_form} leaves exactly one honest emitted notice key: \nolinkurl{worked_example_receipt_interlock_public_request_status_notice}, whose first pointer is \nolinkurl{example_public_request_status_envelope.json}. In the worked successor wrapper, the imported pair \texttt{public\_request\_status\_token = complete} and \texttt{carryforward\_token = preserved\_complete} together with selected family \nolinkurl{worked_example_receipt_interlock_successor_preserved_complete_public_request_notice_normal_form} leaves exactly one honest emitted notice key: \nolinkurl{worked_example_receipt_interlock_public_request_carryforward_notice}, whose first pointer is \nolinkurl{example_public_request_carryforward_envelope.json}. Later terse papers may therefore quote either visible sentence directly only while both the imported basis object and the selected family still agree with that emitted wrapper; if either changes, the wrapper itself has gone stale and must be reopened at the narrower owner.\cite{series:workedexample}


\paragraph{A compact first-reopen-owner card.}
Once one maintained visible request-notice shortcut goes stale, later terse papers still need one small crosswalk: which owner regains authority \emph{first}? The visible-wrapper note should answer that directly rather than forcing every later paper to translate one stale outward-facing sentence back into the right status owner, carryforward owner, canonical-family owner, selector owner, response-handoff owner, or request-side terminal cut from prose.

\begin{table}[t]
\centering
\small
\begin{tabularx}{\linewidth}{@{}p{0.29\linewidth}p{0.22\linewidth}X@{}}
\toprule
Failure class & First reopen owner & Why the visible wrapper is no longer honest \\
\midrule
The same source-only visible wrapper slot still exists but the local emitted notice key, exact outward-facing sentence, wrapper kind, or first reader-pointer order changes & Reopen the public-request-notice owner in Synthesis~61. & The live issue is still the exact visible sentence and pointer order shown to readers. Once those local wrapper anchors drift, the first honest owner remains this note rather than a narrower token owner or a wider request-side handoff note. \\
The wrapper still cites a within-release basis but the imported public request status envelope, its paper-level status token, or its covered-family / family-verdict basis changes & Reopen the public-request-status owner in Synthesis~59 and follow its narrower cutover if needed. & Within-release public request notices compress an already-owned paper-level status envelope rather than proving that token themselves. Once that imported within-release basis drifts, the first honest owner is the status-envelope layer that emits it. \\
The wrapper still cites a successor-release basis but the imported public request carryforward envelope, predecessor/successor token pair, carryforward token, or reopened-family floor changes & Reopen the public-request-carryforward owner in Synthesis~60 and follow its narrower cutover if needed. & Successor-facing public request notices compress an already-owned release-pair carryforward verdict rather than deriving it. Once that imported carryforward basis drifts, the first honest owner is the carryforward layer that emits it. \\
The visible sentence is still the live issue but the reusable canonical family, slot map, or sentence normal form behind that wrapper changes & Reopen the canonical notice-family owner in Synthesis~62. & This note owns one emitted wrapper clause, not the reusable family of sentence shapes that can realize it. Once the family or slot discipline drifts, the first honest owner is the narrower normal-form layer. \\
The reusable family is fixed but the token-keyed choice of which canonical family applies changes & Reopen the notice-selection owner in Synthesis~63. & The visible wrapper depends on one already-selected family choice. Once that selector changes, the first honest owner is the family-selection layer rather than this wrapper clause itself. \\
The visible wrapper basis is fixed and the remaining issue is only the follow-up request-side handoff, predecessor packet cut, successor reuse, refresh delta, visible refresh wrapper, successor-facing handoff, or exact successor-facing packet cut & Widen to Synthesis~64--72 according to which downstream request-side object is now live. & At that point the outward-facing wrapper is no longer the live question. The paper should leave this visible-notice note and widen only to the exact downstream request-side owner whose policy now controls. \\
The visible wrapper basis and downstream request-side branch are fixed and the remaining issue is only request-side terminal reuse, one concrete checked-answer/request-tail seam, or abstract fail-closed cutover & Widen to Synthesis~73 for the maintained request-side terminal stop, drop to Synthesis~17 for one concrete checked-answer/request-tail seam, and widen to Synthesis~74--75 only for the abstract paired-tail discipline or abstract cutover judgment. & Once the visible wrapper and downstream request-side branch are already fixed, later terse papers should stop at the maintained request-side terminal or paired cutover note rather than reopen this wrapper note merely to restate an unchanged sentence. \\
\bottomrule
\end{tabularx}
\caption{A compact first-reopen-owner card. This is the smallest local map from one stale public request notice shortcut to the first narrower or wider owner that regains authority.}
\label{tab:first-reopen-owner-card-public-request-notices}
\end{table}

\paragraph{One tiny first-reopen-owner drill.}
If a later terse paper still asks \emph{``what exact source-only request sentence and first pointer did readers actually see?''} while the surrounding request-side tail remains recognizable but the local emitted notice key, outward-facing sentence, or first pointer changes, Table~\ref{tab:first-reopen-owner-card-public-request-notices} sends the paper back to Synthesis~61 itself. If the wrapper is still the same shape but its imported within-release status envelope or paper-level token drifts, the same table reopens Synthesis~59; if instead its imported successor carryforward basis drifts, it reopens Synthesis~60. If the emitted sentence is still in dispute but the reusable canonical family or slot map changes, the same table reopens Synthesis~62; if the family stays fixed but the selector choosing that family changes, it reopens Synthesis~63. If the visible wrapper basis is already fixed and the real question becomes the downstream request-side handoff or packet chain, the same table widens directly to Synthesis~64--72; if that downstream branch is also fixed and only the request-side terminal stop, one concrete checked-answer/request-tail seam, or the abstract paired cutover remains live, it widens to Synthesis~73 for the maintained request-side terminal stop, drops to Synthesis~17 for one concrete seam, and widens to Synthesis~74--75 only for the abstract paired-terminal discipline or abstract cutover judgment instead of pretending that the visible wrapper still owns the remaining judgment.\cite{series:publicrequeststatusenvelopes,series:publicrequestcarryforward,series:publicrequestnoticenormalforms,series:publicrequestnoticeselections,series:publicrequestresponsemenus,series:publicrequestresponsepackets,series:publicrequestresponsepacketcarryforward,series:publicrequestresponsepacketdeltaledgers,series:publicrequestresponsepacketrefreshnotices,series:publicrequestresponsepacketrefreshnoticenormalforms,series:publicrequestresponsepacketrefreshnoticeselections,series:publicrequestresponsepacketrefreshresponsemenus,series:publicrequestresponsepacketrefreshresponsepackets,series:publicrequestresponsepacketrefreshresponsepacketclosureverdicts,series:pairedterminalcitationdiscipline,series:pairedterminalcutoverrule}

\section{Worked example pointer}
The maintained worked bundle now ships \nolinkurl{example_public_request_notice.json}.\cite{series:workedexample}
It contains two notice entries keyed \nolinkurl{worked_example_receipt_interlock_public_request_status_notice} and \nolinkurl{worked_example_receipt_interlock_public_request_carryforward_notice}.
The within-release entry emits the exact sentence ``For this source-only release, the worked note's paper-level public request promise is complete and supporting artifacts remain available on request.'' and keeps the exact pointer order
\nolinkurl{example_public_request_status_envelope.json} $\rightarrow$ \nolinkurl{example_request_fulfillment_certificates.json} $\rightarrow$ \nolinkurl{example_public_request_contracts.json}.
The successor-facing entry emits the exact sentence ``Across the canned successor release, that complete paper-level public request status carries forward unchanged.'' and keeps the exact pointer order
\nolinkurl{example_public_request_carryforward_envelope.json} $\rightarrow$ \nolinkurl{example_public_request_status_envelope.json} $\rightarrow$ \nolinkurl{example_public_request_contracts.json} $\rightarrow$ \nolinkurl{example_request_fulfillment_certificates.json}.
That gives later terse papers and release notes one compact wrapper object to cite when they mean the visible sentence rather than the narrower status or carryforward owner, while keeping \nolinkurl{complete} and \nolinkurl{preserved_complete} as imported basis vocabulary from Synthesis~59--60 rather than notice-owned semantics.
In the maintained worked bundle, the same within-release and successor-carryforward notice keys are also imported directly onto the shipped audience/question routes so a later terse paper can recover the exact visible wrapper without side-jumping out of the route object. Later terse papers should usually combine those route-visible notice keys with Synthesis~73 for the stable request-side stop, drop to Synthesis~17 for one concrete checked-answer/request-tail seam, and widen to Synthesis~74--75 only for the abstract stop-versus-widen judgment unless the exact sentence or pointer order is itself the dispute.
The companion canonical sentence family for those wrappers is then owned one note later by Synthesis~62, while the token-keyed choice of which such family applies is owned by Synthesis~63.\cite{series:publicrequestnoticenormalforms,series:publicrequestnoticeselections,series:publicrequestresponsepacketrefreshresponsepacketclosureverdicts,series:pairedterminalcitationdiscipline,series:pairedterminalcutoverrule}

\begin{thebibliography}{99}\small

\bibitem{series:publicrequestcontracts}
Anonymous.
\newblock Public Request Contracts for Source-Only Releases: Exact Paper-Level Handoffs from Boilerplate Policy to Maintained Packet Cuts.
\newblock Series synthesis note (Synthesis~56), 2026. (This archive.)

\bibitem{series:requestfulfillmentcertificates}
Anonymous.
\newblock Request Fulfillment Certificates for Source-Only Releases: Exact Completion Criteria for Public Request Contracts and Typed Omitted-Support Cuts.
\newblock Series synthesis note (Synthesis~58), 2026. (This archive.)

\bibitem{series:publicrequeststatusenvelopes}
Anonymous.
\newblock Public Request Status Envelopes for Source-Only Releases: Paper-Level Tokens Compressing Family Completion without Re-owning the Promise.
\newblock Series synthesis note (Synthesis~59), 2026. (This archive.)

\bibitem{series:publicrequestcarryforward}
Anonymous.
\newblock Public Request Carryforward for Source-Only Successor Releases: When Paper-Level Request Status Persists, Advances, or Reopens across Release Evolution.
\newblock Series synthesis note (Synthesis~60), 2026. (This archive.)

\bibitem{series:publicrequestnoticenormalforms}
Anonymous.
\newblock Public Request Notice Normal Forms for Source-Only Releases: Canonical Sentence Families for Within-Release Status and Cross-Release Carryforward Notices.
\newblock Series synthesis note (Synthesis~62), 2026. (This archive.)

\bibitem{series:publicrequestnoticeselections}
Anonymous.
\newblock Public Request Notice Selections for Source-Only Releases: Choosing Canonical Notice Families from Within-Release Status and Cross-Release Carryforward Tokens.
\newblock Series synthesis note (Synthesis~63), 2026. (This archive.)

\bibitem{series:publicrequestresponsemenus}
Anonymous.
\newblock Public Request Response Menus for Source-Only Releases: Smallest-Sufficient Follow-Up Handoffs from Visible Notices to Maintained Request Artifacts.
\newblock Series synthesis note (Synthesis~64), 2026. (This archive.)

\bibitem{series:publicrequestresponsepackets}
Anonymous.
\newblock Public Request Response Packets for Source-Only Releases: Exact Packet Cuts for Maintained Follow-Up Handoffs.
\newblock Series synthesis note (Synthesis~65), 2026. (This archive.)

\bibitem{series:publicrequestresponsepacketcarryforward}
Anonymous.
\newblock Public Request Response Packet Carryforward for Source-Only Successor Releases: When Exact Follow-Up Packets Persist, Refresh, or Reopen.
\newblock Series synthesis note (Synthesis~66), 2026. (This archive.)

\bibitem{series:publicrequestresponsepacketdeltaledgers}
Anonymous.
\newblock Public Request Response Packet Delta Ledgers for Source-Only Successor Releases: Exact Packet-Local Refresh Differences for Maintained Follow-Up Slices.
\newblock Series synthesis note (Synthesis~67), 2026. (This archive.)

\bibitem{series:publicrequestresponsepacketrefreshnotices}
Anonymous.
\newblock Public Request Response Packet Refresh Notices for Source-Only Successor Releases: Visible Reissue Sentences for Refreshed Follow-Up Packets.
\newblock Series synthesis note (Synthesis~68), 2026. (This archive.)

\bibitem{series:publicrequestresponsepacketrefreshnoticenormalforms}
Anonymous.
\newblock Public Request Response Packet Refresh Notice Normal Forms for Source-Only Successor Releases: Canonical Visible Reissue Families for Refreshed Follow-Up Packets.
\newblock Series synthesis note (Synthesis~69), 2026. (This archive.)

\bibitem{series:publicrequestresponsepacketrefreshnoticeselections}
Anonymous.
\newblock Public Request Response Packet Refresh Notice Selections for Source-Only Successor Releases: Choosing Canonical Visible Reissue Families after Packet Refresh.
\newblock Series synthesis note (Synthesis~70), 2026. (This archive.)

\bibitem{series:publicrequestresponsepacketrefreshresponsemenus}
Anonymous.
\newblock Public Request Response Packet Refresh Response Menus for Source-Only Successor Releases: Smallest-Sufficient Successor Handoffs after Visible Refresh.
\newblock Series synthesis note (Synthesis~71), 2026. (This archive.)

\bibitem{series:publicrequestresponsepacketrefreshresponsepackets}
Anonymous.
\newblock Public Request Response Packet Refresh Response Packets for Source-Only Successor Releases: Exact Successor-Facing Packet Cuts after Visible Refresh.
\newblock Series synthesis note (Synthesis~72), 2026. (This archive.)

\bibitem{series:publicrequestresponsepacketrefreshresponsepacketclosureverdicts}
Anonymous.
\newblock Public Request Response Packet Refresh Response Packet Closure Verdicts for Source-Only Successor Releases: Capstone Certificates for Exact Successor-Facing Handoff Slices.
\newblock Series synthesis note (Synthesis~73), 2026. (This archive.)

\bibitem{series:pairedterminalcitationdiscipline}
Anonymous.
\newblock Paired Terminal Citation Discipline for Stable Review Spines: One Answer-Side Stop and One Request-Side Capstone for Terse Papers.
\newblock Series synthesis note (Synthesis~74), 2026. (This archive.)

\bibitem{series:pairedterminalcutoverrule}
Anonymous.
\newblock Paired Terminal Cutover Rules for Stable Review Spines: When Later Terse Papers Must Stop, Widen, or Reopen.
\newblock Series synthesis note (Synthesis~75), 2026. (This archive.)

\bibitem{series:workedexample}
Anonymous.
\newblock Worked Example: Receipt Interlock for an Anonymous-DHT Reader-Privacy Claim.
\newblock Series synthesis note (Synthesis~17), 2026. (This archive.)

\end{thebibliography}
\end{document}
